\section{Instancias, datos y caso base}
\label{sec:instancias}

Los siguientes \textbf{conjuntos de benchmark} se eligen porque aparecen en el paper replicado o
son pruebas mínimas necesarias para validar la implementación. No son el eje teórico del trabajo:
son datos reconocidos que permiten comparar con óptimos publicados y con valores reportados por
la literatura.

\subsection{Benchmarks utilizados}
\begin{description}[leftmargin=1.5em]
  \item[\texttt{toy1} (caso base mínimo).] $N=M=4$, $p=2$, coordenadas 2D. Óptimo a mano $=6$.
        Sirve para validar toda la cadena y derivar los cortes a mano
        (Sección~\ref{sub:verif}).
  \item[OR-Library (Beasley, 1990).] Instancias \texttt{pmed1}--\texttt{pmed40} sobre grafos:
        $N=M\in\{100,200,300,400,500,600,700,800,900\}$, con distintos valores de $p$. Se dan
        como listas de aristas con costo; las distancias $d_{ij}$ se obtienen como \emph{caminos
        mínimos} (Floyd--Warshall). Hay un \textbf{óptimo oficial} por instancia
        (\texttt{pmedopt.txt}) que sirve de referencia dura. En este informe se reporta evidencia
        consolidada: \texttt{pmed1}--\texttt{pmed15}, \texttt{pmed16} como smoke verificado, y
        \texttt{pmed17}--\texttt{pmed40} como extensión local con 24/24 óptimos oficiales.
  \item[TSP-Library (Reinelt, 1991).] \texttt{rl1304} ($N=M=1304$) y \texttt{kroA100} ($N=M=100$),
        con coordenadas 2D y distancia euclidiana \emph{truncada}. Banco principal de escalamiento;
        se resuelve para $p\in\{5,50,200,500\}$ (y $\{10,20,100,300,400\}$ en la comparación con
        el paper).
  \item[RW local mínimo.] \texttt{rw12}: matriz de distancias \textbf{aleatoria y posiblemente
        asimétrica} ($N=M=12$). Valida el caso \emph{no euclidiano / asimétrico} de la
        implementación. No corresponde a la campaña RW grande del paper.
  \item[Grilla sintética euclidiana.] Instancias generadas localmente con seed fijo
        \texttt{20260622}, $N\in\{100,250,500,1000,2000,5000\}$ y
        $p\in\{5\%,10\%,20\%\}$ de $N$. Sirven como prueba de estrés local bajo timeout externo
        de 300 segundos. No son instancias del paper ni tienen óptimos externos conocidos.
\end{description}

\subsection{Supuestos y decisiones de datos}
\begin{itemize}[noitemsep]
  \item \textbf{Distancia euclidiana truncada (floor).} El paper usa $\lfloor\cdot\rfloor$, no el
        redondeo \texttt{nint} canónico de TSPLIB. Se validó: \texttt{rl1304} con $p=5$ y floor da
        $3\,099\,073$, exactamente el óptimo del paper. Se adopta floor en todas las instancias
        con coordenadas por consistencia.
  \item \textbf{Aristas duplicadas en OR-Library.} \texttt{pmed1} trae aristas paralelas con costo
        distinto. La decisión de datos, incorporada aquí para que el informe sea autocontenido, es:
        \emph{la última ocurrencia sobrescribe a las anteriores}. Esta convención coincide con el
        formato de Beasley y con la lectura secuencial del archivo; no se debe reemplazar por
        ``tomar el mínimo''. Si se toma el mínimo, \texttt{pmed1} entrega $5718$ y deja de coincidir
        con el óptimo oficial; con sobrescritura por última ocurrencia entrega $5819$, que sí
        coincide. La implementación verifica además Dijkstra~$\equiv$~Floyd--Warshall en el
        parser. Esta es la decisión que antes estaba documentada como ADR; queda resumida en el
        informe para no depender de anexos externos.
  \item \textbf{Formato propio \texttt{.pmp}.} El solver recibe una primera línea \texttt{N M p}.
        Luego acepta dos variantes: coordenadas \texttt{id x y}, o una matriz explícita precedida
        por la línea \texttt{MATRIX}. Esto cubre instancias geométricas TSPLIB y matrices de
        distancias OR-Library/RW.
\end{itemize}

\subsection{Alcance honesto}
Se ejecutaron: (i) una campaña Benders de 26 casos locales; (ii) evidencia consolidada
OR-Library \texttt{pmed1}--\texttt{pmed40}; (iii) \texttt{rl1304} hasta $p=500$; (iv)
\texttt{kroA100}; (v) una validación asimétrica \texttt{rw12}; y (vi) una grilla sintética local
hasta $N=5000$. \textbf{No} se ejecutaron BIRCH, ODM, RW grande ni TSP enormes del paper; tampoco
se reejecutó Zebra. Por lo tanto no se reportan números locales para esos benchmarks. Cuando se
mencionan, se marcan como resultados del paper o como trabajo futuro.
